\reportheading{3.\quad Task 1 --- Ceiling-jet validation and doorway effects}
\subhead{3.1\quad Agreement with the measured temperature profiles}
At 47.2~kW, the coarse-grid OP-1 profile underestimated the measured temperature rise by 192~K on average. With only R1 open, the signed error averaged +6~K but the RMSE remained 33~K. At 106.2~kW the largest discrepancy occurred under OP-6. The adiabatic calculation placed the hottest station 0.80~m left of the burner. The experiment placed it 0.20~m left. The model captured the direction of displacement and exaggerated its magnitude.
\newcommand{\validationpanel}[3]{%
\begin{tikzpicture}
\begin{axis}[width=.475\linewidth,height=3.9cm,xmin=-340,xmax=340,ymin=0,ymax=1350,
xtick={-300,0,300},ytick={0,600,1200},axis lines=left,ymajorgrids,grid style={gray!14},
tick align=outside,title={#2\quad #3 kW},title style={font=\sffamily\small\bfseries,color=ink,at={(.5,.97)}},
tick label style={font=\sffamily\scriptsize},enlarge x limits=false]
\addplot[teal,thick,no marks] table[col sep=comma,x expr=\thisrow{model_x_cm}-330,y=fds_rise_K]{../../05_Calculations/#1_90_120.csv};
\addplot[only marks,mark=o,mark size=1.2pt,rust] table[col sep=comma,x=figure_r_cm,y=exp_rise_K]{../../05_Calculations/#1_90_120.csv};
\end{axis}\end{tikzpicture}}
\reportfigure{{\sffamily\small Temperature rise (K)\hfill
\textcolor{teal}{\rule[.4ex]{14pt}{1.1pt}} FDS\quad
\textcolor{rust}{$\circ$} Peng et al.\par}
\begin{tabular}{@{}cc@{}}
\validationpanel{op1_47_coarse}{OP-1}{47.2}&\validationpanel{op1_106_coarse}{OP-1}{106.2}\\
\validationpanel{op6_47_coarse}{OP-6}{47.2}&\validationpanel{op6_106_coarse}{OP-6}{106.2}
\end{tabular}\par{\sffamily\small Distance from the fire axis (cm)}}{Figure 1. Measured and simulated ceiling temperature rise. Each panel uses the same axes. The experimental markers are digitised from Peng et al.\ (2021), Fig.~9. The lines show coarse-grid FDS averages over 90--120~s.}

Peng et al.\ reported stable-phase temperatures and illustrated repeatability with OP-1 at 70.8~kW (their Fig.~2). Their reported relative uncertainty of up to 8\% concerns measured temperature. For the four coarse-grid comparisons, we checked the sensitivity to time averaging by comparing 90--105~s with 105--120~s at each station. The largest absolute change at any station ranged from 9.2 to 28.2~K across the four cases. The small signed mean error for OP-6 at low power should therefore be read alongside its spatial RMSE and the choice of averaging window.
\subhead{3.2\quad Sensitivity to grid and model inputs}
For OP-1, coarse, medium and fine grids gave RMSEs of 192.4, 187.7 and 176.2~K at low power, and 82.0, 75.7 and 79.9~K at high power. The high-power error reached its minimum on the medium mesh. Grid refinement alone therefore offers limited grounds for selecting the fine-grid solution. On the high-power coarse grid, switching from Deardorff to WALE reduced RMSE from 82.0 to 77.0~K. Holding Deardorff fixed and changing the radiative fraction from 0.29 to 0.20 or 0.50 gave 74.6 or 131.1~K. Radiative partitioning influenced the temperature prediction more strongly than the tested change of closure in this case.

\clearpage
\subhead{3.3\quad Door exchange and enclosure heat loss}
At 106.2~kW total doorway inflow fell from 0.318~kg/s in OP-1 to 0.095~kg/s in OP-6. At R1, velocity changed from inward to outward between 0.30 and 0.38~m in OP-1 and between 0.22 and 0.30~m in OP-6 (Figure~B1). At 47.2~kW the outward stream at 0.38~m was hotter in OP-6 (212.1 versus 64.8$\,^{\circ}$C). Air entering from R1 could deflect the plume towards the closed side, consistent with the observed peak displacement. The probes establish the doorway exchange; plume deflection is inferred.

For OP-6 at 106.2~kW, the 15~mm proxy board reduced profile RMSE from 361.4 to 100.7~K on the matched diagnostic grid. Peak-position error fell from 0.70 to 0.10~m. The 24 stations from 1.0 to 3.3~m on the R1 side still averaged 112.5~K below the measurements (Figure~2 and Table~F1). Over 90--120~s the board run gave signed domain radiative, convective and conductive terms of $-61.4$, $-34.2$ and $-9.6$~kW, including the exterior extension.
\reportfigure{\begin{tikzpicture}
\begin{axis}[width=.94\linewidth,height=5.7cm,xmin=-340,xmax=340,ymin=0,ymax=1350,ytick={0,300,600,900,1200},xtick={-300,-150,0,150,300},
xlabel={Distance from the fire axis (cm)},ylabel={Temperature rise (K)},axis lines=left,ymajorgrids,grid style={gray!15},
legend columns=3,legend style={at={(.5,1.03)},anchor=south,draw=none,/tikz/every even column/.append style={column sep=6pt}}]
\addplot[gray,thick,dashed] table[col sep=comma,x expr=\thisrow{model_x_cm}-330,y=fds_rise_K]{../../05_Calculations/op6_106_doorflow_90_120.csv};
\addlegendentry{Adiabatic}
\addplot[teal,thick] table[col sep=comma,x expr=\thisrow{model_x_cm}-330,y=fds_rise_K]{../../05_Calculations/op6_106_board_doorflow_90_120.csv};
\addlegendentry{Conducting board}
\addplot[only marks,rust,mark=o,mark size=1.6pt] table[col sep=comma,x=figure_r_cm,y=exp_rise_K]{../../05_Calculations/op6_106_board_doorflow_90_120.csv};
\addlegendentry{Peng et al.}
\end{axis}\end{tikzpicture}}{Figure 2. OP-6 at 106.2 kW. The adiabatic and conducting-board predictions share the doorway diagnostic grid. The measured and simulated profiles use the same coordinate origin as Figure 1.}

On the 47.2~kW OP-1 coarse grid, the same board increased RMSE from 192.4 to 243.3~K. The peak remained over the burner, but its rise fell from 435 to 340~K against a measured 580~K. In the adiabatic OP-1 diagnostic, the 51-station mean cold bias was 191.8~K for the thermocouples and 187.1~K for the co-located gas. The simulated sensor response accounts for only 4.7~K of that difference (Figure~F1). Doorway sensitivity motivated testing both door states and the far ends in Task~2. The propane benchmark assesses ceiling temperature. The propylene clearance margins come from FDS with the prescribed adiabatic enclosure and extraction settings.
